% Copyright (c) 2026 The ALF project.
% This is a part of the ALF project documentation.
% The ALF project documentation by the ALF contributors is licensed
% under a Creative Commons Attribution-ShareAlike 4.0 International License.
% For the licensing details of the documentation see license.CCBYSA.
% !TEX root = doc.tex
%-------------------------------------------------------------------------------------
\section{Delayed and submatrix sequential updates (\texttt{master\_ai})} \label{sec:lowrank}
%-------------------------------------------------------------------------------------
%
\emph{This feature exists only on the AI branches (\texttt{master\_ai}, \texttt{master\_ai\_cpu}). It was written with AI assistance
and is not part of the ALF baseline.}

\paragraph{Purpose.}
The standard sequential update (Sec.~\ref{sec:updating}) applies every accepted local move to the
Green function at once, a rank-$r$ update of cost $N_{\text{dim}}^2$ per move with BLAS level-2
efficiency. The delayed and submatrix schemes keep up to \texttt{N\_delay} accepted moves pending
in low-rank form and apply them together in one matrix product. This changes only the order of the
arithmetic: the acceptance ratios and the Green function are the same quantities, and the random
numbers are drawn in the same order, so a run reproduces the Markov chain of the standard update
up to roundoff.

\paragraph{Build.}
\texttt{source configure.sh <MACHINE> <MODE> LOWRANK}. Without \texttt{LOWRANK} nothing of the
feature is compiled, and the program is identical to the baseline.

\paragraph{Parameters.}
Optional namelist in the file \texttt{parameters}; without it the standard update is used:
\begin{verbatim}
&VAR_LOWRANK
Sequential_scheme = "delayed"   ! "standard" | "delayed" | "submatrix"
N_delay           = 0           ! 0: automatic, -1: automatic and timed, >= 1: fixed
/
\end{verbatim}
The choice is written to the \texttt{info} file. With \texttt{N\_delay = 0} the bound $n_0$ is set from
$N_{\text{dim}}$ and from the largest pending rank one time slice can produce (a fixed rule, the same on
every machine). The best \texttt{N\_delay} depends on the machine (BLAS library, cache sizes), so
\texttt{N\_delay = -1} starts from $n_0$ and measures: in the second and third sweep of the run the time
slices cycle through the candidates $n_0/2$, $n_0$ and $2n_0$ (each limited as the static rule limits
$n_0$), and the wall time per proposed update is summed for each candidate. After the third sweep the
fastest candidate is kept for the rest of the run if it is more than 5\,\% faster than $n_0$; otherwise
$n_0$ is kept. With MPI the times are summed over the ranks of a group, rank 0 of the group decides, and
the choice is broadcast. The result is written to the \texttt{info} file, for example
\texttt{Auto N\_delay: candidates 8 16 32: 1.97E-06 2.21E-06 2.36E-06 s/proposal, chosen 8}.
\texttt{N\_delay} changes only the summation order, so the timing does not change the results beyond
roundoff. A run with fewer than three sweeps ends before the decision and logs nothing; if only one
candidate exists (very few sequential vertices) no timing is done.

\paragraph{Supported cases.}
A vertex is handled by the low-rank scheme if its operator acts on at most 8 sites in every flavour;
other vertices first apply the pending moves and then use the standard update. Vertices updated by
sequential MALA (\texttt{Sequential\_MALA}) are treated the same way: the pending moves are applied,
then the MALA update runs exactly as in the baseline. All pending moves are applied at the end of the
sequential part of every time slice, so the Global\_tau moves (\texttt{Global\_tau\_moves},
\texttt{Global\_tau\_MALA\_moves}), the measurements, the stabilisation and all other parts of ALF see
the full Green function. Global moves, Langevin and HMC do not use the sequential update and are
unaffected. Detailed derivations: \texttt{Documentation/lowrank/*.md}.

\paragraph{Numerics compared with the baseline.}
\begin{itemize}
\item Same quantities, different summation order (roundoff-level differences).
\item Submatrix scheme: blocks with non-diagonal content (bond operators) are applied immediately,
because the submatrix form lost accuracy there: during development, a Kondo run without this rule
reached a Green function precision of $3\times10^{-9}$, against $4\times10^{-11}$ with the standard
update. The delayed scheme has no such loss and is the better choice for models with bond operators.
\end{itemize}

\paragraph{Verification.}
Four levels, from the strictest to the statistical one.
\begin{enumerate}
\item \emph{Bit-identical to the baseline} (\texttt{testsuite/baseline\_guard.py}): a build without the
switch, and a \texttt{LOWRANK} build with \texttt{Sequential\_scheme = "standard"}, give output
bit-identical to \texttt{master}.
\item \emph{Same Markov chain} (\texttt{testsuite/Lowrank.tests/compare\_schemes.py}): with the same seeds the
files \texttt{confout\_*} (final fields and random number generator state) of the delayed and the submatrix
run are byte-identical to the standard run. The proposals of the fields do not depend on the Green
function and all schemes draw the same random numbers in the same order, so identical files mean that
every accept/reject decision was identical. The exception are sequential MALA and Global\_tau MALA, whose proposals contain the force, a
function of the Green function: there roundoff enters the fields and only level 3 is required.
\item \emph{Observables to roundoff}: every number in every bin file agrees with the standard run to
$10^{-8}$ relative (largest value found on all scenarios: $2\times10^{-11}$, usually $10^{-13}$); the
schemes sum in a different order. Unit test \texttt{36-lowrank-unit}: the low-rank kernels against an
exactly recomputed Green function on 23 operator types (diagonal and bond operators, field types
1--4, ill-conditioned cases).
\item \emph{Statistics with independent seeds}: energy and a correlation function of standard and
delayed runs agree within their error bars (separate study, not part of CTest).
\end{enumerate}
\texttt{compare\_schemes.py} runs every scenario with N\_delay automatic and N\_delay $=3$ (many
flushes). The scenarios cover every Hamiltonian of \texttt{Prog/Hamiltonians.list} and the ALF
features of the table below; they are registered in CTest and in the CI of the AI branches for all
compilers (GNU, PGI, Intel, IntelLLVM; serial and MPI).

\paragraph{Compatibility.}
Checked with levels 2 and 3 (GNU/MPICH on a Mac, 2 bins per scenario, 2026-11-20); the final configuration was
identical in all runs except where noted. The largest relative difference of the observables was
$3.5\times10^{-11}$ over 264 scenario/scheme runs of the serial set. In all shipped models every vertex acts on at most
4 sites, so $100\%$ of the vertices use the low-rank kernel; the fallback to the standard update for vertices of more than
8 sites is tested with the test-only model \texttt{LR\_Wide\_Test} (10-site vertices next to one-site vertices, 50\% low rank).
\begin{center}\small
\begin{tabular}{p{0.34\textwidth}p{0.6\textwidth}}
\hline
Hamiltonian & Variants checked \\
\hline
Hubbard & Mz (Ising, type 1), N\_FL $=1$ and 2 (type 2), continuous HS (type 3), N\_SUN 2 and 4, twist and flux
(complex Green function), pinning field; Square, N\_leg\_ladder, Bilayer, Honeycomb, Bilayer honeycomb,
Triangular; projector; Checkerboard and Symm on and off \\
Hubbard\_Plain\_Vanilla & half filling (flavour reconstruction), doped, attractive, adiabatic switching
(time-dependent vertices), projector \\
Kondo & N\_SUN 2 and 4, Symm on and off, $U_c \ne 0$ \\
tV & N\_SUN 1 and 2, attractive and repulsive, bilayer \\
LRC & long-range term on, projector, bilayer, Propose\_S0 \\
Z2\_Matter & gauge fields only, gauge and matter fields, with Hubbard term, projector, Propose\_S0 \\
Spin\_Peierls & square lattice, chain (type 2 and 4 vertices) \\
test models & \texttt{LR\_Jz\_Test} (two-site Jz vertices), \texttt{LR\_Wide\_Test}, \texttt{Hubbard\_Mixed}
(continuous and discrete fields in one run, with sequential MALA) \\
\hline
ALF feature & Result \\
\hline
Global\_tau\_moves (alone and combined with a sequential range), Propose\_S0 & identical chain \\
Global\_moves (implemented only by Hubbard with continuous fields) & identical chain \\
Sequential\_MALA, Global\_tau\_MALA, also with discrete fields in the same run & MALA vertices use the standard
update; the proposal contains the force, so the final configuration agrees only to roundoff
($\le 4\times10^{-11}$) \\
Langevin, HMC & no sequential update, unaffected \\
Restart from \texttt{confin} & identical chain; 2 + 2 bins are byte-identical to 4 bins \\
MPI (2 ranks), TEMPERING (2 replicas, 1 and 2 ranks each), PARALLEL\_PARAMS (own parameter files, also
different models) & identical chain \\
HDF5 & compared in the CI with \texttt{h5dump} \\
\hline
\end{tabular}
\end{center}
Two limits are not of the feature: ALF master itself stops with ``unstable'' when \texttt{Nt\_sequential\_start}
is larger than 1 together with Global\_tau moves in the Hubbard model, and the repulsive tV model ($V=+1$, N\_SUN$=1$) needs
\texttt{Nwrap}$=2$ to reach a Green function precision better than $10^{-6}$ (with \texttt{Nwrap}$=10$ the standard update
itself has a mean precision of $1.7\times10^{-6}$, and the schemes differ by $10^{-6}$ in the observables).
Time-dependent hopping (\texttt{Op\_T} with \texttt{g\_t}) is not used by any shipped model; the wrap of the hopping is the baseline
code in both paths.
The unit test also checks the automatic \texttt{N\_delay} (static rule, timing state machine, error exits), and
\texttt{compare\_schemes.py --n-delay -1} repeats the comparison with the timed automatic \texttt{N\_delay}.
